\section{Capítulo~\ref{sec:acet}. Añadimos \nt{ACET}}

Las tres acciones de la acetilcolina están medidas en rebanadas y en el cerebro vivo; el conflicto entre escritura y lectura se muestra con el modelo del libro; y que \nt{ACET} sirva de línea de habilitación de escritura y comunique la incertidumbre esperada es una hipótesis con apoyo experimental parcial.

{\footnotesize
\setlength{\tabcolsep}{3pt}\renewcommand{\arraystretch}{1.15}
\begin{longtable}{>{\raggedright}p{0.5cm}>{\raggedright}p{4.0cm}>{\raggedright}p{5.6cm}>{\raggedright}p{2.3cm}>{\raggedright\arraybackslash}p{1.7cm}}
\toprule
N.º & Afirmación & Experimento y qué se midió & Fuente & Estado \\
\midrule\endfirsthead
\toprule
N.º & Afirmación & Experimento y qué se midió & Fuente & Estado \\
\midrule\endhead
1 & Las neuronas \nt{ACET} del cerebro son pocas, están reunidas en unos cuantos grupos, y sus salidas se extienden por toda la corteza: un canal de difusión & Tinción con anticuerpos contra la enzima de síntesis de la acetilcolina y marcaje retrógrado en la rata: la corteza, el hipocampo, la amígdala, el bulbo olfatorio y el tálamo reciben acetilcolina sobre todo de seis grupos compactos de células del prosencéfalo basal y del tronco & \cite{Mesulam1983} & medido \\
2 & El nivel de acetilcolina en la corteza es alto en la vigilia y bajo en el sueño lento & Microdiálisis en la corteza y el hipocampo de gatos en libertad de movimiento con registro de EEG: la liberación de acetilcolina en vigilia tranquila es $\sim100\%$ mayor, y en vigilia activa $\sim175\%$ mayor, que en el sueño lento; en el sueño REM, en la corteza, como en vigilia & \cite{Marrosu1995,Hasselmo1999} & medido \\
3 & El sentido de la señal lo fija el receptor: la acetilcolina tiene receptores-canal rápidos y receptores lentos que disparan reacciones en la célula (proposición~\ref{prop:receptor}) & Revisión de medidas: el efecto de la acetilcolina sobre la excitabilidad, la transmisión y la plasticidad depende del lugar de liberación, del subtipo de receptor (los nicotínicos son canales iónicos; los muscarínicos actúan a través de proteínas G) y de la célula destinataria & \cite{Picciotto2012} & medido \\
4 & Primera acción: debilitar las conexiones internas casi sin tocar las entradas externas (factor $1-a$ en~\eqref{eq:acetnet}) & Rebanadas de corteza olfativa de rata: con $100$\,\textmu M de agonistas de la acetilcolina, los potenciales de las conexiones internas (capa Ib) caen más de un $60\%$, y los de la entrada externa (capa Ia), menos de un $15\%$, en la misma célula; la supresión es presináptica. Rebanadas de hipocampo (CA1): $89\%$ frente a $40\%$ para las vías interna y de entrada & \cite{HasselmoBower1992,HasselmoSchnell1994} & medido \\
5 & Segunda acción: reforzar la entrada (factor $\frac{1+a}{2}$) & Rebanadas de neocorteza: los receptores nicotínicos refuerzan las sinapsis de la entrada talámica y no actúan sobre las intracorticales; los muscarínicos debilitan ambos tipos. La acetilcolina aumenta la excitabilidad de las neuronas (revisión) & \cite{Gil1997,Hasselmo2006} & medido \\
6 & Tercera acción: facilitar el cambio de los pesos (tasa $\eta a$) & Rata: la estimulación eléctrica del núcleo basal (la fuente de acetilcolina de la corteza) emparejada con un tono provoca en el animal adulto una fuerte reorganización del mapa de la corteza auditiva hacia el sonido presentado & \cite{KilgardMerzenich1998,Hasselmo2006} & medido \\
7 & La regla de Hebb para patrones dispersos de la segunda ecuación~\eqref{eq:acetnet} da una memoria asociativa que completa el patrón a partir de una parte & Cálculo de la capacidad de una red con patrones dispersos y regla de covarianza: con una fracción pequeña $p$ de neuronas activas, la red guarda bastantes más patrones que con $p=1/2$ & \cite{Tsodyks1988} & teoría \\
8 & Escribir con $a$ bajo daña la memoria: la red escribe lo recordado, no lo visto; con $a=0.1$ el daño no es menor que con $a=0$ & \texttt{models/\allowbreak acet\_memory.py}: $200$ neuronas, cinco patrones de $20$, el patrón nuevo comparte $10$ neuronas con uno de ellos, $10$ ejecuciones. Con $a=0$ el patrón nuevo no se memoriza, y el viejo arrastra $5.9$ neuronas ajenas de $10$; con $a=0.1$ el nuevo se recuerda ($0.97$), pero ambos se funden: el nuevo arrastra $7.3$ neuronas ajenas y el viejo, $7.9$; desde $a=0.2$, menos de una & deducción en el capítulo & modelo \\
9 & Proposición~\ref{prop:writeread}: no hay un nivel $a$ con el que escritura y lectura sean igual de buenas (fig.~\ref{fig:acetsim}) & El mismo script, tasa de aprendizaje $\eta a$: el patrón nuevo se memoriza entero en una presentación con $a\ge0.9$, en $0.8$ con $a=0.75$, y no se memoriza con $a\le0.5$. Lectura a partir de la mitad del patrón: $1.0$ con $a\le0.2$, $0.96$ con $a=0.3$, $0.04$ con $a=0.4$ & deducción en el capítulo & modelo \\
10 & En el cerebro, un solo cable de difusión, \nt{ACET}, conmuta escritura y lectura & La idea procede de modelos construidos a partir de medidas en rebanadas. El experimento más directo es en humanos: el bloqueante de receptores muscarínicos escopolamina empeora el aprendizaje de pares de palabras nuevos, sobre todo de los que más se solapan con los viejos, pero no el recuerdo de pares ya aprendidos. No hay medida directa de los dos modos de la red & \cite{HasselmoAndersonBower1992,Atri2004} & hipótesis \\
11 & El filtro~\eqref{eq:kalman} es óptimo; el peso de la entrada y la tasa de aprendizaje son una misma magnitud $P/R$; en régimen estacionario $P=\sqrt{qR}$ y la memoria es $\sqrt{R/q}$ (proposición~\ref{prop:kalman}) & No requiere experimento: la optimalidad del filtro de Kalman--Bucy es un resultado clásico de la teoría de estimación; el régimen estacionario se deduce en el capítulo & \cite{KalmanBucy1961} & teoría \\
12 & \nt{ACET} comunica a la red una magnitud parecida a $P$: la incertidumbre esperada & Propuesto en un modelo probabilístico de la atención. No hay medida directa de que el nivel de acetilcolina siga la incertidumbre de la estimación & \cite{YuDayan2005} & hipótesis \\
13 & Parte de las neuronas \nt{ACET} responden a la recompensa y al castigo en un par de decenas de milisegundos, con más fuerza cuanto más inesperado es el refuerzo & Ratón, neuronas colinérgicas del prosencéfalo basal identificadas optogenéticamente: respuesta a la recompensa y al castigo a los $18\pm3$\,ms; la magnitud crece con lo inesperado del refuerzo; las respuestas son parecidas en neuronas de dos núcleos & \cite{Hangya2015} & medido \\
14 & \nt{NORA} nota un cambio inesperado del mundo, y \nt{ACET} mantiene la red en modo de escritura & El reparto del trabajo se propone en un modelo. Indirectamente: en humanos, el diámetro de la pupila, ligado a \nt{NORA}, sigue los cambios y la tasa de aprendizaje. No hay comprobación directa del par \nt{NORA}--\nt{ACET} & \cite{YuDayan2005,Nassar2012} & hipótesis \\
15 & En el sueño lento la red, en modo de lectura, reproduce lo escrito durante el día, y estas repeticiones lo copian en la memoria a largo plazo & Registros de conjuntos de neuronas del hipocampo de rata: las células que trabajaron juntas de día se disparan juntas con más frecuencia en el sueño lento posterior, y en el mismo orden: medido. Para la reescritura: suprimir las ondas agudas del hipocampo tras el aprendizaje empeora la memoria espacial, y alargar las ondas espontáneas la mejora; que la causa sea el nivel bajo de \nt{ACET} no está demostrado & \cite{WilsonMcNaughton1994,Skaggs1996,Girardeau2009,FernandezRuiz2019,Hasselmo1999} & hipótesis \\
\bottomrule
\end{longtable}}
